\chapter{Ke Mana Arah AI}
\label{ch:terbaru}

Selama dua tahun kuliah, bidang ini bergerak lebih cepat daripada slide yang mengajarkannya. Ketika
kuliah LSTM membahas xLSTM sebagai pesaing transformer, model bahasa besar sudah dipakai jutaan orang
setiap hari. Ketika kuliah AI and law membahas usulan aturan untuk model serbaguna, aturan itu sudah
berlaku. Bab penutup ini mencoba merangkum perkembangan besar beberapa tahun terakhir dan menghubungkannya
kembali dengan fondasi di bab-bab sebelumnya. Isinya pasti cepat usang, dan memang itulah intinya.

\section{Model fondasi dan hukum penskalaan}

Istilah \istilah{foundation model} \citep{bommasani2021} menamai pergeseran cara membangun AI. Alih-alih
melatih model khusus untuk setiap tugas, kita melatih satu model besar pada data yang luas,
biasanya dengan self-supervised learning (\cref{ch:data}), lalu menyesuaikannya untuk banyak tugas. GPT-4
\citep{openai2023}, model-model terbuka seperti LLaMA \citep{touvron2023}, CLIP untuk gambar dan teks
\citep{radford2021}, dan AlphaFold untuk protein adalah contoh-contohnya.

Pendorong utamanya adalah pengamatan empiris yang disebut \istilah{scaling laws}. \citet{kaplan2020}
menemukan bahwa loss model bahasa turun mengikuti hukum pangkat terhadap jumlah parameter $N$, jumlah
data $D$, dan komputasi $C$, misalnya $L(N)\approx(N_c/N)^{\alpha}$ dengan $\alpha$ kecil. Hukum pangkat
berarti garis lurus pada skala log--log: setiap kali model diperbesar sepuluh kali, loss turun dengan
faktor yang kira-kira tetap. Komputasi pelatihan transformer bisa diperkirakan dengan aturan praktis
$C\approx6ND$ operasi floating point: sekitar dua operasi per parameter per token untuk lintasan maju,
dan sekitar empat untuk lintasan mundur. \citet{hoffmann2022} menunjukkan bahwa untuk anggaran komputasi
tertentu, banyak model sebelumnya terlalu besar dan dilatih dengan terlalu sedikit data, dan bahwa
jumlah token yang optimal kira-kira dua puluh kali jumlah parameter. Model dengan 70 miliar parameter
dengan demikian butuh sekitar 1,4 triliun token, dan komputasinya $6\times7\times10^{10}\times1{,}4\times10^{12}
\approx5{,}9\times10^{23}$ operasi. Angka ini menarik dibandingkan dengan ambang $10^{25}$ operasi yang dipakai
AI Act untuk menandai model dengan risiko sistemik (\cref{ch:manusia}): ambang itu kira-kira tujuh belas
kali lebih besar.

Supaya model raksasa lebih murah dijalankan, arsitektur \istilah{mixture of experts} \citep{fedus2022}
hanya mengaktifkan sebagian kecil subjaringan untuk setiap token, sehingga jumlah parameter total bisa
besar sekali sementara komputasi per token tetap terkendali. Untuk menyesuaikan model ke tugas baru tanpa
melatih ulang semua parameternya, LoRA \citep{hu2022} membekukan matriks bobot $\mW\in\R^{d\times d}$ dan
hanya mempelajari koreksi berpangkat rendah $\bm B\bm A$ dengan $\bm B\in\R^{d\times r}$ dan
$\bm A\in\R^{r\times d}$. Untuk $d=4096$ dan $r=8$, satu matriks penuh punya sekitar 16,8 juta parameter,
sedangkan koreksinya hanya $2\times4096\times8=65\,536$, kurang dari setengah persen.

\section{Menyelaraskan model dengan manusia}

Model bahasa yang hanya dilatih menebak kata berikutnya belum tentu membantu, jujur, atau aman.
InstructGPT \citep{ouyang2022} memperkenalkan resep tiga langkah: penyetelan dengan contoh jawaban yang
ditulis manusia, pelatihan model ganjaran dari perbandingan pasangan jawaban \citep{christiano2017}, lalu
optimasi model bahasa terhadap model ganjaran itu dengan PPO (\cref{ch:rl}), sambil menjaga agar ia tidak
terlalu jauh dari model awal.

Direct preference optimization \citep{rafailov2023} memperlihatkan bahwa langkah reinforcement learning
bisa dilewati. Penurunannya elegan. Tujuan RLHF adalah memaksimalkan ganjaran harapan dikurangi penalti
KL terhadap model acuan $\pi_{\mathrm{ref}}$:
\begin{equation}
\max_\pi\ \E_{y\sim\pi(\cdot\mid x)}\big[r(x,y)\big]-\beta\,\KL\big(\pi(\cdot\mid x)\,\|\,\pi_{\mathrm{ref}}(\cdot\mid x)\big).
\end{equation}
Masalah ini punya solusi tertutup, $\pi^*(y\mid x)\propto\pi_{\mathrm{ref}}(y\mid x)\exp\big(r(x,y)/\beta\big)$.
Membalik persamaan ini memberi ganjaran dalam bentuk policy, $r(x,y)=\beta\log\big(\pi^*(y\mid x)/\pi_{\mathrm{ref}}(y\mid x)\big)$
ditambah suku yang hanya bergantung pada $x$. Dalam model preferensi Bradley--Terry, peluang jawaban $y_w$
lebih disukai dari $y_l$ adalah $\sigma\big(r(x,y_w)-r(x,y_l)\big)$, dan suku yang hanya bergantung pada $x$
saling menghapus dalam selisih itu. Hasilnya, preferensi bisa dipelajari langsung pada policy dengan loss
\begin{equation}
\Loss_{\mathrm{DPO}}=-\log\sigma\Big(\beta\log\frac{\pi_\theta(y_w\mid x)}{\pi_{\mathrm{ref}}(y_w\mid x)}-\beta\log\frac{\pi_\theta(y_l\mid x)}{\pi_{\mathrm{ref}}(y_l\mid x)}\Big),
\end{equation}
sebuah binary cross-entropy biasa, tanpa model ganjaran terpisah dan tanpa sampling selama pelatihan.

\section{Model yang bernalar}

\citet{wei2022} menunjukkan bahwa meminta model menuliskan langkah-langkah penalarannya sebelum menjawab,
\istilah{chain-of-thought}, memperbaiki kinerja pada soal matematika dan logika secara nyata. Generasi
berikutnya melatih kemampuan ini secara langsung. DeepSeek-R1 \citep{deepseek2025} memperlihatkan bahwa
reinforcement learning dengan ganjaran yang bisa diperiksa secara otomatis, misalnya apakah jawaban akhir
sebuah soal matematika benar atau apakah sebuah program lolos uji, membuat model belajar menulis rantai
penalaran panjang, memeriksa ulang, dan mencoba jalan lain ketika buntu. Pola ini mengingatkan pada
pelajaran dari AlphaZero di \cref{ch:klasik}: memberi model lebih banyak waktu untuk ``berpikir'' saat
menjawab, dengan mencari dan memeriksa, menjadi sumbu penskalaan baru di samping ukuran model.

Hasilnya terlihat jelas dalam matematika. AlphaGeometry \citep{trinh2024} menggabungkan model bahasa dengan
mesin deduksi simbolik untuk menyelesaikan soal geometri olimpiade. Pada Olimpiade Matematika
Internasional Juli 2025, sebuah versi Gemini Deep Think dari Google DeepMind menyelesaikan lima dari
enam soal dalam bahasa alami dan diakui panitia mencapai standar medali emas\footnote{Pengumuman Google
DeepMind, Juli 2025: \url{https://deepmind.google/blog/advanced-version-of-gemini-with-deep-think-officially-achieves-gold-medal-standard-at-the-international-mathematical-olympiad/}.
OpenAI melaporkan hasil setara dengan model eksperimentalnya, tetapi hasil itu tidak disertifikasi oleh
koordinator IMO.}.

\section{Agen}

Model bahasa yang bisa memakai alat berubah dari penjawab pertanyaan menjadi agen. ReAct \citep{yao2023}
merangkai penalaran dan tindakan secara bergantian: model menulis pikiran, memanggil alat seperti
pencarian web atau kalkulator, membaca hasilnya, lalu berpikir lagi. Retrieval-augmented generation
\citep{lewis2020} memberi model akses ke basis dokumen sehingga jawabannya bisa bersandar pada sumber yang
dikutip, bukan hanya pada ingatan parameternya. Agen pemrograman yang bisa membaca kode, menjalankan uji,
dan memperbaiki kesalahannya sendiri menjadi salah satu penerapan yang paling cepat berkembang. Sebagian
besar masalah lama muncul kembali dalam bentuk baru: credit assignment untuk tugas panjang (\cref{ch:rl}),
galat yang menumpuk sepanjang banyak langkah seperti pada rollout simulator (\cref{ch:operator}), dan
pertanyaan siapa yang bertanggung jawab ketika agen bertindak (\cref{ch:manusia}).

\section{Gambar, video, dan model dunia}

Model difusi dari \cref{ch:mlat} menjadi mesin di balik generator gambar modern. Latent diffusion
\citep{rombach2022} menjalankan difusi di ruang laten sebuah autoencoder, bukan di ruang piksel, sehingga
jauh lebih murah. Diffusion transformer \citep{peebles2023} mengganti U-Net dengan transformer dan
memperlihatkan bahwa model difusi juga mengikuti hukum penskalaan. Generator video adalah perpanjangan
alami dari gagasan yang sama ke dimensi waktu.

Model yang membangkitkan video yang konsisten secara fisis mulai disebut \istilah{world models}, sebuah
istilah yang dipakai \citet{ha2018} untuk agen yang belajar model internal lingkungannya lalu berlatih di
dalam ``mimpi''-nya sendiri. Apakah model video benar-benar memahami fisika atau hanya meniru
penampilannya adalah pertanyaan terbuka, dan alat dari Bagian V, seperti uji kekekalan dan ekstrapolasi,
adalah cara yang baik untuk mengujinya.

\section{AI untuk sains}

Tahun 2024 memberi penanda simbolis. Hadiah Nobel Fisika diberikan kepada John Hopfield dan Geoffrey
Hinton untuk penemuan dasar yang memungkinkan machine learning dengan jaringan saraf tiruan, dan
setengah Hadiah Nobel Kimia diberikan kepada Demis Hassabis dan John Jumper untuk prediksi struktur
protein (\cref{ch:protein}). Jaringan Hopfield yang dihargai itu adalah leluhur langsung jaringan Hopfield
modern di \cref{ch:lstm}.

Di bidang cuaca, model berbasis machine learning seperti GraphCast \citep{lam2023}, GenCast
\citep{price2025}, dan Aurora \citep{bodnar2024} kini bersaing dengan sistem numerik operasional, dan
beberapa pusat prakiraan cuaca mulai menjalankannya berdampingan. Model fondasi untuk fisika umum
\citep{mccabe2023,alkin2024} dan kumpulan data simulasi besar \citep{ohana2024} mencoba mengulang
kesuksesan itu untuk PDE secara umum. Di kimia dan biologi, AlphaFold3 \citep{abramson2024} memperluas
prediksi struktur ke interaksi antar molekul. Pola yang sama muncul di semua bidang ini: data simulasi
dan pengukuran yang terkumpul selama puluhan tahun menjadi bahan bakar, dan struktur fisika, seperti
simetri dan kekekalan, menjadi kemudi.

\section{Yang tidak berubah}

Di tengah semua perubahan ini, pelajaran dari kuliah-kuliah dasar tetap berlaku, bahkan semakin penting.
Generalisasi di luar data latih tetap sulit (\cref{ch:data}). Model besar tetap bisa belajar jalan pintas
yang salah, dan evaluasi yang jujur tetap memerlukan pembagian data yang hati-hati. Optimasi tetap
bergantung pada condition number dan noise (\cref{ch:optimasi}). Model generatif tetap bisa
berhalusinasi, sehingga penalaran yang bisa diperiksa dan sumber yang bisa dikutip menjadi semakin
berharga. Dan sistem yang kuat tetap harus dijelaskan dengan jujur kepada orang-orang yang hidupnya
terdampak (\cref{ch:xai,ch:manusia}).

Kalau ada satu benang merah yang saya bawa dari empat semester ini, mungkin begini: kemajuan terbesar
jarang datang dari membuang apa yang sudah diketahui. LSTM bertahan karena constant error carousel adalah
gagasan yang benar. Konvolusi bertahan karena simetri geser adalah sifat dunia. Persamaan Navier--Stokes
tidak berhenti berlaku karena ada jaringan saraf. Model terbaik adalah yang belajar dari data sambil
menghormati apa yang sudah kita ketahui tentang dunia.

\sumberbab{Bab ini disusun dari bacaan, bukan dari satu mata kuliah. Sebagian topiknya disinggung di kuliah
LSTM (xLSTM dan Mamba), NLP (model bahasa besar), AI and Law (model AI serbaguna), dan AI in Life Sciences
(model bahasa untuk kimia). Perkembangan yang baru diperiksa dengan pengumuman resminya pada saat
penulisan, Oktober 2026.}
